\ifdefined\MyWebBook\else
\documentclass[../textbook.tex]{subfiles}
\begin{document}
\fi
\bookappendix{模型资源估算}{app:resource-estimation}{2}

本附录汇总参数量、显存与计算时间的估算公式及预算记录项。估算结果取决于模型结构、执行阶段、数值格式、批量与序列长度分布；理论存储量、运行时峰值与实测值分别反映静态规模、执行需求和硬件表现。公式推导见\bookxref{ch:resource-budget}，硬件供给、并行切分和服务测量分别见\bookxref{ch:compute-infrastructure}、\bookxref{ch:distributed-training}与\bookxref{ch:serving-architecture}。

\section{计量约定}

\begin{table}[htbp]
\centering
\caption{资源估算的计量对象、单位与换算关系}\label{tab:appendix-resource-units}
\begin{tabularx}{\linewidth}{p{.22\linewidth}p{.28\linewidth}X}
\toprule
对象 & 常用单位 & 含义与换算\\\midrule
参数数量 & 个参数，常写 M 或 B & $1\mathrm{M}=10^6$，$1\mathrm{B}=10^9$；参数数本身不是字节数。\\
存储量 & B、MB、GB、MiB、GiB & $1\mathrm{GB}=10^9$ B，$1\mathrm{GiB}=2^{30}$ B；报告应保留原始字节数或明确换算。\\
计算工作量 & FLOPs & 一次乘法和一次加法按本书口径计为两个浮点操作。\\
计算速率 & FLOP/s & 完成工作量的速度；必须注明数值精度、稀疏性及峰值或实测口径。\\
带宽 & B/s、GB/s & 单位时间可搬运的字节数；显存、互连、网络和存储带宽属于不同层级。\\
延迟 & s、ms & 需声明起止边界、预热、并发、输入输出长度及分位数定义。\\
吞吐 & token/s、request/s & 单请求、单设备和集群总吞吐不可混写；输入与输出词元可分别报告。\\
\bottomrule
\end{tabularx}
\end{table}

若一种数据类型或状态每元素占 $b$ 字节、共有 $N$ 个元素，则理想连续存储为
\begin{equation}
 M=Nb\quad\text{字节}.
 \label{eq:appendix-elements-bytes}
\end{equation}
实际分配还可能包含对齐、填充、量化元数据、内存碎片、框架保留池和临时工作区。估算表应把式\eqref{eq:appendix-elements-bytes}得到的理论值与运行时附加项分列。

\section{参数量计算}

考虑层数 $L$、隐藏宽度 $d$、词表大小 $V$、查询头数 $H_q$、键值头数 $H_{kv}$、每头宽度 $d_h$ 和前馈宽度 $d_f$ 的自回归 Transformer。若 $H_qd_h=d$，单层注意力参数为
\begin{equation}
 P_{\mathrm{attn}}=2d^2+2dH_{kv}d_h.
 \label{eq:appendix-attn-params}
\end{equation}
前两项中的一个 $d^2$ 来自查询投影、另一个来自输出投影，键和值各贡献 $dH_{kv}d_h$。标准 MHA 取 $H_{kv}=H_q$，此时注意力合计约 $4d^2$；GQA 和 MQA 只减少键值投影部分。

设 $c_f=2$ 表示普通两矩阵前馈，$c_f=3$ 表示两条上投影加一条下投影的门控前馈，则
\begin{equation}
 P_{\mathrm{ffn}}=c_fdd_f.
\end{equation}
输入嵌入与输出词表投影共享时记 $\tau=1$，独立时记 $\tau=2$。忽略偏置并把归一化、位置参数和任务头单列，可写成
\begin{equation}
 P\approx \tau Vd+L\left(2d^2+2dH_{kv}d_h+c_fdd_f\right)
 +P_{\mathrm{norm}}+P_{\mathrm{extra}}.
 \label{eq:appendix-total-params}
\end{equation}

参数量取决于输入输出权重共享、前馈结构、键值头数和跨层共享方式。偏置、查询键归一化、位置表、专家、适配器与任务头的参数按实际结构分别计入。MoE 的专家总参数量由专家数 $E$ 决定，每词元参与计算的专家参数量则由激活数 $k$ 决定。

\begin{table}[htbp]
\centering
\caption{常见模块的参数量}\label{tab:appendix-parameter-rules}
\begin{tabularx}{\linewidth}{p{.25\linewidth}p{.27\linewidth}X}
\toprule
模块 & 参数量 & 适用条件\\\midrule
线性层 $a\rightarrow b$ & $ab$；有偏置时加 $b$ & 参数存储方向不改变独立标量数量。\\
嵌入表 & $Vd$ & 查表计算量不按一次完整 $Vd$ 矩阵乘法计算。\\
普通前馈 & $2dd_f$ & 一次上投影、一次下投影。\\
门控前馈 & $3dd_f$ & 两条上投影、一次下投影。\\
LoRA 分支 & $r(d_{\mathrm{in}}+d_{\mathrm{out}})$ & 每个目标矩阵单独计算；缩放系数通常不是大项。\\
门控专家集合 & $E\,3dd_e$ & 每个专家采用门控前馈且参数互不共享；共享专家与路由器另计。\\
ViT 补丁投影 & $(CP_hP_w)d+d$ & 含偏置；与补丁数量 $N$ 无关。\\
\bottomrule
\end{tabularx}
\end{table}

\section{模型状态的显存开销}

推理权重的基础存储为
\begin{equation}
 M_w=Pb_w.
\end{equation}
若采用 $q$ 位分组量化、每 $g$ 个参数共享一个占 $b_s$ 字节的尺度，则近似为
\begin{equation}
 M_w\approx\frac{Pq}{8}+\frac{P}{g}b_s+M_{\mathrm{zero}}+M_{\mathrm{unquantized}}+M_{\mathrm{packing}}.
 \label{eq:appendix-quant-memory}
\end{equation}
零点、未量化层、打包与对齐必须按实际格式加入；有些后端还在执行时建立反量化工作区。

训练状态应从实际保存的数组求和：
\begin{equation}
 M_{\mathrm{state}}=P_wb_w+P_gb_g+\sum_jP_{o,j}b_{o,j}.
 \label{eq:appendix-training-state}
\end{equation}
这里 $P_w,P_g,P_{o,j}$ 分别是权重、梯度与第 $j$ 个优化器状态的元素数。全参数训练时它们通常与 $P$ 同阶；冻结参数、分片参数和量化基座需要分别计数。

\begin{table}[htbp]
\centering
\caption{混合精度 Adam 每参数状态示例}\label{tab:appendix-adam-bytes}
\begin{tabular}{lcc}
\toprule
状态 & 每元素字节 & 每参数贡献\\\midrule
计算权重 & 2 & 2 B\\
梯度 & 2 & 2 B\\
高精度主权重 & 4 & 4 B\\
一阶矩 & 4 & 4 B\\
二阶矩 & 4 & 4 B\\\midrule
合计 & — & 16 B\\
\bottomrule
\end{tabular}
\end{table}

表\ref{tab:appendix-adam-bytes}只描述一种配置。若梯度以 FP32 累积、没有独立主权重、使用低精度优化器或采用不同分片策略，每参数字节数都会变化。LoRA/QLoRA 仍需保留基座权重，并为穿过冻结模块的反向路径保存或重算激活，不能把全参数训练显存简单乘以可训练参数比例。

在理想均匀分片条件下，各数据并行方案的单设备常驻状态存储量为
\begin{align}
 M_{\mathrm{DDP}}&=M_w+M_g+M_o,\\
 M_{\mathrm{Z1}}&=M_w+M_g+\frac{M_o}{N_d},\\
 M_{\mathrm{Z2}}&=M_w+\frac{M_g+M_o}{N_d},\\
 M_{\mathrm{Z3}}&=\frac{M_w+M_g+M_o}{N_d},
\end{align}
其中 $N_d$ 是数据并行组大小。单设备峰值还需加入参数聚合、通信桶、预取、激活和工作区；这些项不应统一除以设备总数。

\section{运行时显存开销}

训练激活存储量的线性近似为
\begin{equation}
 M_{\mathrm{act}}\approx c_aLBTdb_a
 \label{eq:appendix-activation-memory}
\end{equation}
其中，系数 $c_a$ 表示每层每位置实际保存的隐藏宽度等效元素数，其取值由模型结构、融合算子和重计算策略决定，并通过实测峰值校准。若显式保存注意力概率，还可能出现 $O(LBH_qT^2)$ 元素；分块注意力和激活重计算会改变保存集合。

同构 GQA 的 KV 缓存对第 $i$ 个请求已处理长度 $T_i$ 的理论字节数为
\begin{equation}
 M_{\mathrm{KV}}=2LH_{kv}d_hb_c\sum_{i=1}^{B}T_i.
 \label{eq:appendix-kv-memory}
\end{equation}
系数2对应键和值。所有请求等长时才可化为 $2LBTH_{kv}d_hb_c$。分页缓存还要把每条请求长度向块容量取整；共享前缀按物理页计一次，不能按逻辑引用数重复相加。

训练峰值可组织为
\begin{equation}
 M_{\mathrm{train,peak}}=
 M_{\mathrm{resident}}+M_{\mathrm{act}}+M_{\mathrm{gather}}
 +M_{\mathrm{comm}}+M_{\mathrm{workspace}}+M_{\mathrm{reserve}}.
 \label{eq:appendix-training-peak}
\end{equation}
推理峰值通常包含权重、KV 缓存、当前批次激活、采样或束搜索状态、工作区和运行时保留。各项峰值之和给出保守估计，实际峰值取决于这些状态在执行过程中的同时驻留情况。

\section{计算时间估算}

矩阵 $[n,a]$ 乘 $[a,b]$ 的前向工作量约为
\begin{equation}
 F_{\mathrm{matmul}}\approx2nab.
\end{equation}
对应线性层的输入梯度和权重梯度各有一次同阶矩阵乘法，所以前向加反向约为 $6nab$。稠密 Transformer 训练常以
\begin{equation}
 F_{\mathrm{train}}\approx6PD
 \label{eq:appendix-six-pd}
\end{equation}
作为矩阵计算基线，其中 $D$ 是处理词元数。该式未自动包含长序列注意力交互、激活重算、路由、Softmax、逐元素运算和优化器更新。

常规密集注意力每层前向的两个交互矩阵乘法近似为
\begin{equation}
 F_{\mathrm{interaction}}\approx4BT^2d.
 \label{eq:appendix-attention-flops}
\end{equation}
长上下文下应把它与线性投影工作量分别列出。因果掩码是否减少实际算术取决于内核是否跳过未来区域，仅在分数计算后写入负无穷不会省去矩阵乘法。

已知模型工作量 $F$、设备数 $N$、单设备对应精度峰值速率 $R_{\mathrm{peak}}$ 和按同一工作量口径定义的有效利用率 $u$，墙钟时间估算为
\begin{equation}
 t\approx\frac{F}{NR_{\mathrm{peak}}u}.
 \label{eq:appendix-runtime}
\end{equation}
若 $u$ 已由端到端墙钟时间反推，它已经吸收通信、等待和运行时开销；再次乘同一类折损系数会重复扣减。扩展设备数后，$u$ 通常随拓扑、批量和并行方案变化。

带宽下界可与计算下界并列：若必须搬运 $Q$ 字节、有效带宽为 $\beta$，则
\begin{equation}
 t\geq\max\left(\frac{F}{R_{\mathrm{effective}}},\frac{Q}{\beta}\right).
\end{equation}
预填充通常具有较高矩阵复用，逐词元解码则常反复读取权重与缓存。两阶段应分别测量首词元时间、逐输出词元时间和吞吐。

\begin{example}[权重、训练状态与 KV 缓存的容量估算]

设模型有 $P=7\times10^9$ 个参数。两字节权重的理论存储为
\begin{equation*}
 7\times10^9\times2=\SI{14e9}{\byte}\approx\SI{13.04}{\gibi\byte}.
\end{equation*}
若采用表\ref{tab:appendix-adam-bytes}的16字节训练状态配置，则未分片状态约为
\begin{equation*}
 \SI{112e9}{\byte}\approx\SI{104.31}{\gibi\byte}.
\end{equation*}
这两个结果都没有包含激活、通信缓冲、工作区和运行时保留。

再设 $L=32,H_{kv}=8,d_h=128$，KV 缓存采用两字节元素。每请求每位置为
\begin{equation*}
 2(32)(8)(128)(2)=\SI{131072}{\byte}=\SI{128}{\kibi\byte}.
\end{equation*}
单请求缓存8192个位置正好为1 GiB；16个同长请求为16 GiB。分页取整、共享前缀和分布式切分会改变物理驻留量，因此部署容量还要代入实际长度分布和缓存管理方式。

\end{example}

\section{资源预算记录表}

\begin{longtable}{@{}p{.25\linewidth}p{.27\linewidth}p{.39\linewidth}@{}}
\caption{资源预算的最小记录字段}\label{tab:appendix-resource-worksheet}\\
\toprule
项目 & 填写值 & 口径与核对问题\\\midrule
\endfirsthead
\toprule
项目 & 填写值 & 口径与核对问题\\\midrule
\endhead
模型身份 & 名称、版本、配置哈希 & 实际层数、宽度、词表、头数、前馈和共享关系是否与名称对应？\\
执行阶段 & 训练、预填充、解码或评测 & 各阶段是否分别预算，是否包含反向与优化器更新？\\
输入负载 & $B,T$ 或长度分布 & 填充、有效词元、并发、输出上限和长尾场景如何定义？\\
数值格式 & 权重、梯度、状态、激活、缓存 & 位数之外是否存在尺度、零点、主权重及未量化层？\\
参数与权重 & $P$；字节、GiB & 独立参数与执行参数是否分开，输入输出是否共享？\\
训练状态 & 分项字节数 & 梯度、主权重和各优化器状态是否按实际精度计数？\\
激活与缓存 & 公式值、校准系数 & 是否采用重计算、分块注意力、分页或前缀共享？\\
临时与保留 & 聚合、通信、工作区、余量 & 是否记录了峰值同时性和框架内存保留？\\
计算工作量 & FLOPs 与包含项 & 是否包含长度平方项、重计算和其他额外执行？\\
硬件供给 & 数量、容量、峰值率、带宽、拓扑 & 峰值率是否对应实际精度，跨设备路径是否明确？\\
时间估算 & $u$、估计秒数、期限 & 利用率口径从何而来，扩容后是否重新估计？\\
实测校准 & 峰值显存、吞吐、延迟分位数 & 预热、采样窗口、同步边界和监控工具是否固定？\\
\bottomrule
\end{longtable}

资源预算的结论包括容量约束、完成时间和待校准系数。参数量描述独立标量的数量，显存需求取决于执行状态的驻留方式，计算时间取决于工作量、数据搬运与设备利用率。

\chapterreferences
\myendchapterdocument
